% ----------------------------------------------------------------
% AMS-LaTeX definitions for a Dissertation/Thesis ****************
% Template prepared by T.A. McWalter (March 2009)
% ----------------------------------------------------------------
\documentclass[a4paper,11pt,oneside]{report}

% PACKAGES -------------------------------------------------------

\usepackage{amsmath}
\usepackage{amsthm}
\usepackage{amsfonts}
\usepackage{amssymb}
\usepackage{fancyhdr}
\usepackage{ifpdf}
\usepackage{mathrsfs}
\usepackage[numbers]{natbib}
\usepackage{xspace}

% GRAPHICS HANDLING ----------------------------------------------

\ifpdf
    \usepackage[pdftex]{graphicx}
    \DeclareGraphicsExtensions{.pdf, .jpg, .tif}
    \newcommand{\bb}{viewport}
\else
    \usepackage{graphicx}
    \DeclareGraphicsExtensions{.eps, .jpg}
    \newcommand{\bb}{bb}
\fi

% OTHER SET-UP ---------------------------------------------------

\pagestyle{fancy}
\lhead{\rightmark}
\chead{}
\rhead{\thepage}
\cfoot{}

\setlength{\topmargin}{0.5cm}
\setlength{\headsep}{1cm}
\setlength{\textwidth}{14.0cm}
\setlength{\headwidth}{14.0cm}
\setlength{\headheight}{14pt}
\setlength{\textheight}{22cm}
\setlength{\oddsidemargin}{1.3cm}
\setlength{\evensidemargin}{0cm}

\newcommand{\setlinespacing}[1]
           {\renewcommand{\baselinestretch}{#1}\small\normalsize}

% THEOREM STYLES -------------------------------------------------

\theoremstyle{plain}
\newtheorem{thm}{Theorem}[chapter]
\newtheorem{prop}[thm]{Proposition}
\newtheorem{cor}[thm]{Corollary}
\newtheorem{con}[thm]{Conjecture}
\newtheorem{lem}[thm]{Lemma}
\theoremstyle{definition}
\newtheorem{defn}[thm]{Definition}
\newtheorem{ass}[thm]{Assumption}
\newtheorem{exam}[thm]{Example}

% MACRO DEFINITIONS ----------------------------------------------

% Statistical
\def\E{\mathbb{E}}
\def\Var{\mathbb{V}\mathtt{ar}}

% Derivatives
\def\S{\ensuremath{S_t}\xspace}
\def\dS{\ensuremath{dS_t}\xspace}
\def\dt{\ensuremath{dt}\xspace}
\def\dW{\ensuremath{dW_t}\xspace}
\def\dq{\ensuremath{dq_t}\xspace}

% Differences
\def\DS{\ensuremath{\delta S_t}\xspace}
\def\Dt{\ensuremath{\delta t}\xspace}
\def\DW{\ensuremath{\delta W_t}\xspace}
\def\Dq{\ensuremath{\delta q_t}\xspace}

% Binomial model
\def\u{\ensuremath{u}\xspace}
\def\d{\ensuremath{d}\xspace}
\def\o{\ensuremath{o}\xspace}
\def\pu{\ensuremath{p_\u}\xspace}
\def\pd{\ensuremath{p_\d}\xspace}
\def\po{\ensuremath{p_\o}\xspace}

% PREAMBLE -------------------------------------------------------

\title{Valuation of Convertible Bonds with Credit Risk}

\author{David Naylor\\[2cm]
Programme in Advanced Mathematics of Finance,\\
University of the Witwatersrand,\\
Johannesburg.\\[2cm]
\includegraphics[width=4cm]{../common/WitsColourLogo}\\[2cm]
}

% MAIN DOCUMENT --------------------------------------------------

\begin{document}

\maketitle \setlinespacing{1.20}

% DECLARATION ----------------------------------------------------

\def\abstractname{\Huge Declaration}
\begin{abstract}
I declare that this project is my own, unaided work. It is being submitted as partial fulfilment of the Degree of Bachelor of Science with Honours in the University of the Witwatersrand, Johannesburg. It has not been submitted before for any degree or examination in any other University.\\[3cm]
\noindent\rule{5cm}{0.5pt}\\[1cm]
\noindent \today
\end{abstract}

% ABSTRACT -------------------------------------------------------

\def\abstractname{\Huge Abstract}
\begin{abstract}
Place abstract here.  Based on paper \cite{AFV03}.
\end{abstract}

% ACKNOWLEDGEMENTS -----------------------------------------------

\def\abstractname{\Huge Acknowledgements}
\begin{abstract}
Place acknowledgements here.
\end{abstract}

% TABLE OF CONTENTS ----------------------------------------------

\setlinespacing{1}
\pagenumbering{roman}
\tableofcontents
\listoffigures
\listoftables
\clearpage
\setlinespacing{1.20}
\pagenumbering{arabic}
\setcounter{page}{1}

% CONTENT --------------------------------------------------------

\chapter{Binomial Model with Credit Risk}
\section{Specification}
Consider a stock with the following stochastic model\footnote{Of note $\mathbb{E}[\frac{\dS}{\S}] = r\dt \Leftrightarrow \frac{\dS}{\S} = (r + \lambda\eta)\dt + \sigma\dW - \eta\dq$}:
\begin{displaymath}
 \frac{\dS}{\S} = (r + \lambda\eta)\dt + \sigma\dW - \eta\dq
\end{displaymath}
where:

\begin{tabular}{ll}
 $r$            & is the risk free rate \\
 $\sigma$       & is the log-volatility of the stock price \\
 $\lambda$      & is the hazard rate \\
 $\eta$         & is the percentage drop in stock price on a default event \\
 \dW            & is a Wiener process \\
 \dq            & is a Poisson jump process where the first jump is the default event \\
\end{tabular}

\noindent and
\begin{displaymath}\begin{array}{rcl}
 \E[\frac{\dS}{\S}]     &=& r\dt \\
 \E[\frac{S_{t + \dt}}{\S}] &=& \E[\frac{\S}{\S} + \frac{\dS}{\S}] \\
                        &=& 1 + \E[\frac{\dS}{\S}] \\
 \Var[\frac{\dS}{\S}]   &=& (\sigma^2 + \lambda\eta^2)\dt \\
 \Var[\frac{S_{t + \dt}}{\S}] &=& \Var[\frac{\S}{\S} + \frac{\dS}{\S}] \\
                        &=& \Var[\frac{\dS}{\S}] \\
\end{array}\end{displaymath}

\noindent and consider a binomial model that has time step \Dt, jumps \u and \d with probability \pu and \pd\footnote{$\pd = 1 - \pu - \po = e^{-\lambda\Dt} - \pu$} respectively and where the probability of default is \po\footnote{The time till the first jump follows an exponential distribution and has probability $1 - e^{-\lambda\Dt}$} and (TODO: insert picture illustration)
\begin{displaymath}\begin{array}{rcl}
 \E[\frac{S_{t + \Dt}}{\S}]     &=& \u\pu + \d\pd + (1 - \eta)\po \\
                                &=& \u\pu + \d(e^{-\lambda\Dt} - \pu) + (1 - \eta)(1 - e^{-\lambda\Dt}) \\
 \Var[\frac{S_{t + \Dt}}{\S}]   &=& \u^2\pu + \d^2\pd + (1 - \eta)^2\po - \E^2[\frac{\DS}{\S}] \\
                                &=& \u^2\pu + \d^2(e^{-\lambda\Dt} - \pu) + (1 - \eta)^2(1 - e^{-\lambda\Dt}) - \E^2[\frac{\DS}{\S}] \\
\end{array}\end{displaymath}

\section{Derivation}
If we equate the first moment of the stochastic model with that of the binomial model, in \Dt time, we get:
\begin{displaymath}\begin{array}{rcl}
 e^{r\Dt}       &=& \u\pu + \d(e^{-\lambda\Dt} - \pu) + (1 - \eta)(1 - e^{-\lambda\Dt}) \\
 \pu(\u - \d)   &=& e^{r\Dt} - \d e^{-\lambda\Dt} - (1 - \eta)(1 - e^{-\lambda\Dt}) \\
 \pu            &=& \frac{e^{r\Dt} - \d e^{-\lambda\Dt} - (1 - \eta)(1 - e^{-\lambda\Dt})}{\u - \d} \\
 \Rightarrow\pd &=& -\frac{e^{r\Dt} - \u e^{-\lambda\Dt} - (1 - \eta)(1 - e^{-\lambda\Dt})}{\u - \d} \\
\end{array}\end{displaymath}
and if we equate the second moment about the mean of the stochastic model with that of the binomial model, in \Dt time, we get:
\begin{displaymath}\begin{array}{rcl}
 (\sigma^2 + \lambda\eta^2)\Dt &=& \u^2\pu + \d^2(e^{-\lambda\Dt} - \pu) + (1 - \eta)^2(1 - e^{-\lambda\Dt}) - \E^2[\frac{\DS}{\S}] \\
                        &=& (\u^2 - \d^2)\pu + \d^2 e^{-\lambda\Dt} +  (1 - \eta)^2(1 - e^{-\lambda\Dt}) - e^{-2r\Dt} \\
                        &=& (\u + \d)(e^{r\Dt} - \d e^{-\lambda\Dt} - (1 - \eta)(1 - e^{-\lambda\Dt})) \\
                        & & + \d^2 e^{-\lambda\Dt} + (1 - \eta)^2(1 - e^{-\lambda\Dt}) - e^{-2r\Dt} \\
                        &=& (\u + \d)(e^{r\Dt} - (1 - \eta)(1 - e^{-\lambda\Dt})) - \u\d e^{-\lambda\Dt} \\
                        & & + (1 - \eta)^2(1 - e^{-\lambda\Dt}) - e^{-2r\Dt} \\
\end{array}\end{displaymath}
If we make the assumption $\Dt^2 = 0$, $\u\d = 1$ and $\u = e^{\sqrt{A\Dt}}$ and replace all exponentials with Taylor series expansion we get:
\begin{displaymath}\begin{array}{rcl}
 \u                     &=& 1 + \sqrt{A\Dt} + \frac{A\Dt}{2!} + \frac{(A\Dt)^{\frac{3}{2}}}{3!} \\
 \d                     &=& 1 - \sqrt{A\Dt} + \frac{A\Dt}{2!} - \frac{(A\Dt)^{\frac{3}{2}}}{3!} \\
 \u + \d                &=& 2 + A\Dt \\
 (\sigma^2 + \lambda\eta^2)\Dt &=& (2 + A\Dt)(1 + r\Dt - \lambda\Dt(1 - \eta)) - (1 - \lambda\Dt) \\
                        & & + \lambda\Dt(1 - \eta)^2 - (1 - 2r\Dt) \\
                        &=& A\Dt + \lambda\Dt - \lambda\Dt(1 - \eta)(1 + \eta) \\
                        &=& A\Dt + \lambda\eta^2\Dt \\
 A\Dt                  &=& \sigma^2\Dt
\end{array}\end{displaymath}

\section{Parameterisation}
Based on the above the following table specifies a binomial model with credit risk (TODO: include table in List of Tables):

\begin{center}\begin{tabular}{|rcl|l|l|}\hline
 \multicolumn{3}{|l|}{\textbf{Parameter}} & \textbf{Description} \\\hline
 \u     &=& $e^{\sigma\sqrt{\Dt}}$ &
    Multiplier for an up-step \\\hline
 \d     &=& $e^{-\sigma\sqrt{\Dt}}$ &
    Multiplier for a down-step \\\hline
 \pu    &=& $\frac{e^{r\Dt} - \d e^{-\lambda\Dt} - (1 - \eta)(1 - e^{-\lambda\Dt})}{\u - \d}$ &
    Probability of an up-step \\\hline
 \pd    &=& $e^{-\lambda\Dt} - \pu$ &
    Probability of a down-step \\\hline
 \po    &=& $1 - e^{-\lambda\Dt}$ &
    Probability of a default \\\hline
 \Dt    & & &
    Time-step \\ \hline
\end{tabular}\end{center}

This model was first derived by \citet{MK12}, who showed that this model converges, in continuous time, to the \citet{AFV03} model.  \Citet{MK12} also showed the valuation method of this model is the same as in the classical binomial model:
\begin{displaymath}
 V_t = e^{-r\Dt}(\pu V^{+}_t + \pd V^{-}_t- + \po X_t)
\end{displaymath}
where

\begin{tabular}{ll}
 $V_t$          & is the value at time $t$ \\
 $V^+_t$        & is the value of an up-step at time $t + \Dt$ \\
 $V^-_t$        & is the value of a down-step at time $t + \Dt$ \\
 $X_t$          & is the value of default at time $t$ \\
\end{tabular}


% BIBLIOGRAPHY ---------------------------------------------------

\clearpage
\setlinespacing{1}
\addcontentsline{toc}{chapter}{Bibliography}
\bibliographystyle{../references/amsplnat}
\bibliography{../references/references.bib}
\end{document}

%
% Document ends here - everything beyond this point is ignored.
%


% GRAPHIC TEMPLATE FOR EPS OR PDF FILE ---------------------------

\begin{figure}[h!]
    \centering
    \includegraphics[width=8cm]{filename}
    \caption[Short caption for figure list]{Main caption.}\label{FigLabel}
\end{figure}

% GRAPHIC TEMPLATE FOR JPG FILE ----------------------------------

\begin{figure}[h!]
    \centering
    \includegraphics[\bb=0 0 XXX YYY, width=8cm]{filename.jpg}
    \caption[Short caption for figure list]{Main caption.}\label{FigLabel}
\end{figure}

% Replace XXX and YYY with pixel width and height of image.

